\chapter{Fysikaliska konstanter}
\begin{table}[h]
\raggedright
\small
\begin{tabularx}{\textwidth}{p{6.5cm}X}
\toprule
\textbf{Konstant} & \textbf{Värde} \\
\midrule
Absoluta nollpunkten & $T_0 = -273,15$ °C = 0 K \\ %Checked
Universella massaheten för atommassa & $u = 1,66053906892 \cdot 10^{-27}$ kg \\ %Checked
Ljushastigheten i vakuum & $c = 299\,792\,458$ m/s \\ %Checked
Tyngdaccelerationen normalvärde & $g = 9,80665$ m/s$^2$ \\ %Checked
Allmänna gaskonstanten & $R_u = 8,314462618$ J/(mol$\cdot$K) \\ %Checked
Avogadros tal & $N_A = 6,02214076 \cdot 10^{23}$ mol$^{-1}$ \\ %Checked
Molvolymen vid ideal gas vid 0°C och 1 bar & $22,71095464 \cdot 10^{-3}$ m$^3$/mol \\ %Checked
Faradays konstant & $F = 96485,33212$ (A$\cdot$s)/mol eller C/mol \\ %Checked
Boltzmanns konstant & $k = \frac{R_u}{N_A} = 1,380649 \cdot 10^{-23}$ J/K \\
Permeabiliten i vakuum & $\mu_0 = 4\pi \cdot 10^{-7} = 1,25663706127 \cdot 10^{-6}$ (V$\cdot$s)/(A$\cdot$m) \\ %Checked
Dielektricitetskonstanten & $\varepsilon_0 = 8,8541878188 \cdot 10^{-12}$ (A$\cdot$s)/(V$\cdot$m) \\ %Checked
Elektriska elementarladdningen & $e = 1,602176634 \cdot 10^{-19}$ (A$\cdot$s) \\ %Checked
Elektronens vilomassa & $m_e = 9,1093837139 \cdot 10^{-31}$ kg \\ %Checked
Kvoten mellan elektronens laddning och massa & $e/m = 1,75882000838 \cdot 10^{11}$ (A$\cdot$s)/kg \\ %Checked
Neutronens vilomassa & $m_n = 1,67492750056 \cdot 10^{-27}$ kg \\ %Checked
Neutronens vilomassa & $m_n = 1,00866491606$ u \\ %Checked
Protonens vilomassa & $m_p = 1,67262192595 \cdot 10^{-27}$ kg \\ %Checked
Protonens vilomassa & $m_p = 1.0072764665789$ u \\ %Checked
Kvoten mellan protonens och elektronens massa & $m_p/m_e = 1836,152673426$ \\ %Checked
Plancks konstant & $h = 6,62607015 \cdot 10^{-34}$ J$\cdot$s \\ %Checked
Stefan-Boltzmanns konstant & $\sigma = 5,670374419 \cdot 10^{-8}$ W/(m$^2\cdot$K$^4$) \\ %Checked
Total solinstrålning (TSI) & $G_{sc} = 1361,6 \pm 0,3 $ W/m$^2$ \\
\bottomrule
\end{tabularx}
\end{table}
(\cite{NIST_Constants_2026}, \cite{NASA_TSI_2026})
